\documentclass[11pt]{article}
\usepackage[margin=1in]{geometry}
\usepackage{graphicx,booktabs,float,xcolor,hyperref}
\usepackage{helvet}\renewcommand{\familydefault}{\sfdefault}
\hypersetup{colorlinks=true,linkcolor=blue,urlcolor=blue}
\newcommand{\Mw}{$M_{\mathrm w}$}
\setlength{\parindent}{0pt}\setlength{\parskip}{6pt}
\title{\textbf{Prospective source model of the 2026 \Mw~6.8 Kumamoto (Japan) earthquake}\\
\large A blind test of the MTUQ$+$WISP pipeline, produced before the USGS finite fault}
\author{Automated pipeline (17.Venezuela\_2026)}
\date{Frozen 2026-07-28, hours after origin --- for comparison with the USGS solution when posted}
\begin{document}\maketitle

\section*{Purpose}
No USGS finite-fault model existed for \texttt{us6000tgb9} when this was produced. We run the
teleseismic pipeline on open GSN/FDSN data and freeze the prediction so it can be compared with the
USGS (and JMA/NIED) solutions once published. Only teleseismic body $+$ surface waves and the public
W-phase moment tensor (as seed) are used --- no strong motion, no GNSS, no InSAR.

\section{Event and frozen benchmark}
\begin{table}[H]\centering\small
\begin{tabular}{ll}
\toprule
Origin & 2026-07-28T07:27:15.5 UTC, 32.682$^\circ$N, 130.722$^\circ$E, depth 10 km \\
Magnitude & \Mw~6.8 (W-phase); $M_0 = 1.73\times10^{19}$ N\,m; centroid 11.5 km; DC 87.6\,\% \\
Tectonics & shallow crustal \textbf{strike-slip}, Kumamoto fault zone (2016 \Mw 7.0 sequence region) \\
W-phase NP1 & 212.5/75.1/$-$163.8 \quad right-lateral, NE--SW Futagawa--Hinagu trend \\
W-phase NP2 & 118.2/74.3/$-$15.5 \quad left-lateral conjugate \\
\bottomrule
\end{tabular}
\end{table}

\section{Data}
The rewritten fetcher delivered the best coverage of any event in this project: \textbf{241
three-component stations (641 channels), 32--89$^\circ$, maximum azimuthal gap 30$^\circ$}. WISP
applied its own quality control on top; MTUQ used response-corrected subsets of 25 and 60 stations.

\begin{figure}[H]\centering
\includegraphics[width=0.7\textwidth]{../fig_kumamoto_stations.png}
\caption{Teleseismic stations (azimuthal-equidistant; rings 30/60/90$^\circ$).}
\end{figure}

\section{Finite fault (WISP), both nodal planes, \texttt{shift\_match} applied}
WISP \texttt{auto\_model} on both W-phase planes, seeded from the W-phase moment, followed by the
cross-correlation time-shift step and re-inversion (standing rule; it improved the misfit by
9\,\% and 6\,\%).

\begin{table}[H]\centering\small
\begin{tabular}{llllll}
\toprule
plane & strike/dip/rake & \Mw & misfit (pre$\to$post shift) & peak slip & T$_{5-95}$ / V$_r$ \\
\midrule
\textbf{NP2 = Futagawa trend} & \textbf{212/75/$-$164} & 6.75 & 0.1607 $\to$ \textbf{0.1507} & 1.88 m & 13.6 s / 2.05 \\
NP1 = conjugate & 118/74/$-$16 & 6.75 & 0.1667 $\to$ 0.1523 & 1.63 m & 14.8 s / 2.09 \\
\bottomrule
\end{tabular}
\caption{The tectonically expected right-lateral plane fits marginally better ($\sim$1\,\%). As in
our Morocco/Chile/Mexico studies the teleseismic fit is a weak plane discriminant; here the weak
preference at least points the same way as the regional tectonics.}
\end{table}

\begin{figure}[H]\centering
\includegraphics[width=0.88\textwidth]{../20260728072715/ffm.0/NP2_shift/plots/SlipIsochrone_trueScale.png}
\caption{Preferred plane (212/75/$-$164): single compact asperity slightly updip of the 10-km
hypocentre, slip $\sim$2--16 km depth, peak 1.88 m, rupture front spanning $\sim$14 s at
V$_r\approx$2.1 km s$^{-1}$ --- a smaller sibling of the 2016 Kumamoto rupture.}
\end{figure}
\begin{figure}[H]\centering
\includegraphics[width=0.55\textwidth]{../20260728072715/ffm.0/NP2_shift/plots/MomentRate.png}
\caption{Moment-rate function (preferred plane).}
\end{figure}
\begin{figure}[H]\centering
\includegraphics[width=0.92\textwidth]{../20260728072715/ffm.0/NP2_shift/plots/Map.png}
\caption{Map view of the preferred-plane slip.}
\end{figure}
\begin{figure}[H]\centering
\includegraphics[width=0.92\textwidth]{../20260728072715/ffm.0/NP2_shift/plots/P_body_waves.png}
\caption{Teleseismic P fits (preferred plane, shifted).}
\end{figure}
\begin{figure}[H]\centering
\includegraphics[width=0.92\textwidth]{../20260728072715/ffm.0/NP2_shift/plots/SH_body_waves.png}
\caption{Teleseismic SH fits (preferred plane, shifted).}
\end{figure}

\section{Moment tensor (MTUQ) --- an honest failure mode, documented}
The independent MTUQ mechanism estimates were \emph{unstable} for this event, in an instructive way:

\begin{table}[H]\centering\small
\begin{tabular}{lllll}
\toprule
variant & stations / gap & mechanism & \Mw & Kagan vs W-phase \\
\midrule
DC, body P only & 25 / 97$^\circ$ & 94/59/$-$29 & (edge) & 27.1$^\circ$ \\
DC, body P$+$SH & 25 / 97$^\circ$ & 58/50/$-$83 & 6.30 & 67.1$^\circ$ \\
DC, joint & 25 / 97$^\circ$ & 184/82/11 & 6.65 & $\sim$78$^\circ$ \\
full MT (surface) & 25 / 97$^\circ$ & 198/81/9 ($\gamma$=$\delta$=0) & 6.60 & 78.0$^\circ$ \\
\textbf{DC, joint, 60 stations} & \textbf{60 / 80$^\circ$} & \textbf{130/88/7} & \textbf{6.80} & \textbf{27.3$^\circ$} \\
DC, body P$+$SH, 60 stations & 60 / 80$^\circ$ & 212/85/$-$2 & 6.75 & 91.1$^\circ$ \\
\bottomrule
\end{tabular}
\caption{MTUQ variants. With the 25-station subset (14 of 25 stations inside one 45$^\circ$
southern azimuth band) the surface-wave solutions lock onto a mechanism rotated $\sim$90$^\circ$
from the W-phase --- the four-lobe aliasing expected for vertical strike-slip radiation, whose
patterns repeat every 90$^\circ$ of azimuth. Doubling the stations (gap 97$^\circ\to$80$^\circ$) snapped the joint solution out of the rotated minimum (78$^\circ\to$27.3$^\circ$) and recovered $M_w$ 6.80 --- confirming the aliasing diagnosis. The body P$+$SH variant remains aliased; teleseismic SH is unusable at this azimuth geometry. The pipeline's independent mechanism precision here ($\sim$27$^\circ$) is honestly below the 5--19$^\circ$ achieved for the thrust events.}
\end{table}

Two things follow. First, for \textbf{strike-slip} events azimuthal \emph{balance} matters more
than for dip-slip events, because clustered sampling aliases the four-lobed radiation pattern;
station count alone is not protection. Second, the shallow-event ``body-only'' rule from our
Morocco/Mexico work is a \emph{dip-slip} cure (Kanamori--Given indeterminacy of $M_{rt},M_{rp}$)
and does not transfer to vertical strike-slip, where the long-period surface waves are, in
principle, well behaved --- given adequate azimuthal sampling.

\section{Frozen prospective prediction}
\begin{itemize}
\item \textbf{Fault plane:} right-lateral \textbf{212/75/$-$164} (Futagawa--Hinagu trend), weakly
preferred by the waveform fit and strongly by tectonic context.
\item \textbf{\Mw:} 6.75 (WISP, both planes), vs W-phase 6.8.
\item \textbf{Slip:} single compact asperity, peak $\sim$1.9 m, $\sim$2--16 km depth, centred at /
slightly updip of the hypocentre; no shallow slip maximum at the surface.
\item \textbf{Duration:} T$_{5-95}\approx$14 s, moment-rate peak at $\sim$3.5 s; V$_r\approx$2.1
km s$^{-1}$.
\end{itemize}

\textbf{Caveats.} Teleseismic-only; the plane preference is $\sim$1\,\% in misfit (weak); the
independent point-source mechanism could not be stabilised from the azimuthally clustered
long-period subset (documented above), so the mechanism statement leans on the W-phase planes plus
the relative WISP fit. USGS will likely add regional/InSAR data; per our Al Haouz study their model
may be more compact than ours.

\end{document}
